\chapter{Appendix: glossary of escape names}
\label{ch:escglossary}

\begin{aside}
Every escape category and its full name. Decoding cryptic
sequences becomes easier when you know what the letters stand
for.
\end{aside}

\section{Prefix codes}

\begin{itemize}[leftmargin=*]
  \item \code{ESC} = \code{0x1B} (escape).
  \item \code{CSI} = \code{ESC [} (Control Sequence
    Introducer).
  \item \code{OSC} = \code{ESC ]} (Operating System
    Command).
  \item \code{DCS} = \code{ESC P} (Device Control String).
  \item \code{APC} = \code{ESC \_} (Application Program
    Command).
  \item \code{ST} = \code{ESC \textbackslash} (String
    Terminator).
  \item \code{BEL} = \code{0x07} (bell).
\end{itemize}

\section{CSI commands}

\begin{itemize}[leftmargin=*]
  \item \code{A}: cursor up.
  \item \code{B}: cursor down.
  \item \code{C}: cursor forward.
  \item \code{D}: cursor backward.
  \item \code{H}: cursor position.
  \item \code{J}: erase in display.
  \item \code{K}: erase in line.
  \item \code{m}: SGR (style).
  \item \code{n}: device status report.
  \item \code{r}: set scroll region.
  \item \code{s}: save cursor.
  \item \code{u}: restore cursor (or CSI-u keyboard).
\end{itemize}

\section{SGR}

\code{SGR} = Select Graphics Rendition. The escape
\code{ESC [ N m} where N is a parameter selecting
attributes or colours.

\section{DEC private modes}

DEC private modes use \code{ESC [ ? N h/l}. Examples:
\code{?25} cursor visibility, \code{?1049} alt screen,
\code{?2026} sync output.

\section{OSC commands}

\begin{itemize}[leftmargin=*]
  \item \code{0}: window title.
  \item \code{1}: icon title.
  \item \code{2}: window title.
  \item \code{4}: define palette colour.
  \item \code{7}: current working directory.
  \item \code{8}: hyperlink.
  \item \code{9}: notification (iTerm2).
  \item \code{52}: clipboard.
\end{itemize}

\section{Recap}

\begin{recap}
\begin{itemize}[leftmargin=*]
  \item Letter codes have full names.
  \item Knowing them speeds decoding ANSI streams.
\end{itemize}
\end{recap}
